\documentclass[11pt]{article}
\usepackage[margin=1in]{geometry}
\usepackage{amsmath,amssymb}
\usepackage{enumitem}
\setlength{\parindent}{0pt}\setlength{\parskip}{6pt}
\setlist[enumerate,1]{label=\textbf{\arabic*.},leftmargin=*,itemsep=1.1em}
\setlist[enumerate,2]{label=(\alph*),leftmargin=2em,itemsep=0.5em,topsep=0.3em}
\newcommand{\duedate}{Tuesday, September 15, by 11:59PM}
\begin{document}
\begin{center}{\Large\bfseries Math 465: Homework 2}\\[4pt]{\large Due: \duedate}\end{center}
\hrule\vspace{0.8em}
\textbf{Instructions.} Write full and complete solutions to the following problems and submit them on
Gradescope by the deadline. For each problem, explain the reasoning behind your
answer; a formula or numerical answer without justification is not a complete
solution. \textbf{If your final answer is numerical, please carry out the
arithmetic} and state the resulting number (so, for example, if your answer is
$6!-5!$, I want you to also tell me the answer is $600$.)

You may discuss the problems with other students, but you must write up your
solutions independently and in your own words. You are not allowed to use
generative AI tools, Internet searches, or any other external means to directly
obtain a solution. Generative AI tools may also not be used to proofread your
writing or \LaTeX. If you are unsure about what is or is not permitted, please ask.

At the beginning of your submission, acknowledge any permitted books, notes, or
other sources you consulted, and list the names of any students with whom you
discussed the homework.
\vspace{0.5em}\hrule\vspace{1em}
\begin{enumerate}
\item Recall that a (strong) composition of $n$ is a sequence $a=(a_1,\ldots,a_k)$ of positive integers such that $a_1+\cdots+a_k=n$. The entries $a_i$ are called parts.
\begin{enumerate}
\item[(a)] How many total compositions does $n$ have?\\
  There are $n-1$ spaces between $n$ consecutive units. Choosing which spaces contain separators gives $2^{n-1}$ compositions.
\item[(b)] For $n\ge 2$, how many compositions of $n$ are there with an even number of even parts?\\
  Give each composition weight $(-1)^r$, where $r$ is its number of even parts. The signed generating function is
\[
\sum_{n\ge0}(E_n-O_n)x^n=\frac{1}{1-(x+x^3+\cdots)+(x^2+x^4+\cdots)}=\frac{1-x^2}{1-x}=1+x.
\]
Thus $E_n-O_n=0$ for $n\ge2$. Since $E_n+O_n=2^{n-1}$, we get $E_n=2^{n-2}$.
\item[(c)] Find a bijection between the set of compositions of $n$ with an even number of even parts and the set of binary strings of length $n-2$.\\
  Encode a composition by its separator word $b_1\cdots b_{n-1}$ and delete $b_1$. For each fixed $b_2\cdots b_{n-1}$, the two choices of $b_1$ split or join the first block, so exactly one choice gives an even number of even parts. The remaining $n-2$ bits give a binary string, and the construction is reversible.
\end{enumerate}
\item Show that the number of weak compositions of $n$ with $k$ parts equals the number of weak compositions of $k-1$ with $n+1$ parts.\\
  By stars and bars, the two numbers are respectively
\[
\binom{n+k-1}{k-1}\quad\text{and}\quad \binom{(k-1)+(n+1)-1}{(n+1)-1}=\binom{n+k-1}{n}.
\]
These are equal by symmetry of binomial coefficients.
\item How many binary strings of length 12 have four $0$'s and eight $1$'s with no consecutive $0$'s?\\
  Arrange the eight 1's first. They create nine gaps, including the two ends. Choose four distinct gaps for the zeros, giving $\binom94=126$.
\item How many lattice paths from $(0,0)$ to $(8,5)$, using only steps one unit east and one unit north, pass through the point $(3,2)?$\\
  Split the path at $(3,2)$. The number is
\[
\binom{5}{3}\binom{8}{5}=10\cdot56=560.
\]
\item A lattice path from $(0,0)$ to $(m,n)$ uses only steps one unit east and one unit north. Assume $m,n\ge2$. How many such paths have exactly three changes of direction?\\
  There are four alternating runs. Each starting direction contributes $(m-1)(n-1)$ choices for the two positive east and two positive north run lengths. Therefore the answer is $2(m-1)(n-1)$.
\item Prove the identity
\[
\sum_{j=0}^{n}\binom{m+j}{m}=\binom{m+n+1}{m+1}
\]
by counting lattice paths two ways.\\
  Count paths from $(0,0)$ to $(m+1,n)$ by the height $j$ of the last north step. The preceding path has $m$ east and $j$ north steps, giving $\binom{m+j}{m}$ choices. Directly, the number of paths is $\binom{m+n+1}{m+1}$.
\item Let $C_n$ be the number of north-east lattice paths from $(0,0)$ to $(n,n)$ that never go above $x=y$. Prove that
\[
C_n=\binom{2n}{n}-\binom{2n}{n-1}=\frac{1}{n+1}\binom{2n}{n}.
\]\\
  There are $\binom{2n}{n}$ total paths. Reflecting the initial part through the diagonal up to the first step above it bijects bad paths with paths to $(n-1,n+1)$, counted by $\binom{2n}{n-1}$. Hence
\[
C_n=\binom{2n}{n}-\binom{2n}{n-1}=\frac{1}{n+1}\binom{2n}{n}.
\]
\item Evaluate
\[
\sum_{k=0}^{n}\binom{n}{k}2^k
\]
and give a combinatorial explanation.\\
  The binomial theorem gives $(1+2)^n=3^n$. Each object is omitted or selected with one of two labels, and $\binom nk2^k$ counts outcomes with $k$ selected objects.
\item Find the coefficient of $x^5$ in the generalized binomial expansion of
\[
(1+x)^{-3/2}.
\]\\
  By the generalized binomial theorem,
\[
[x^5](1+x)^{-3/2}=\binom{-3/2}{5}=-\frac{63}{256}.
\]
\item Prove
\[
\frac{1}{\sqrt{1-4x}}=\sum_{n\ge0}\binom{2n}{n}x^n.
\]\\
  The generalized binomial theorem gives $(1-4x)^{-1/2}=\sum_{n\ge0}\binom{-1/2}{n}(-4x)^n$. Since $\binom{-1/2}{n}(-4)^n=\frac{(2n)!}{(n!)^2}=\binom{2n}{n}$, the identity follows.
\end{enumerate}
\end{document}
